\documentclass{article}
\usepackage[T1]{fontenc}
\usepackage{lmodern}
\usepackage{graphicx}
\usepackage{amsmath,amssymb}
\usepackage[a4paper,margin=2cm]{geometry}
\usepackage[absolute,overlay]{textpos}
\setlength{\TPHorizModule}{1mm}
\setlength{\TPVertModule}{1mm}
\usepackage[indonesian]{babel}
\usepackage{hyperref}
\setlength{\emergencystretch}{2em}
% unit: Halaman 59

\begin{document}

% === Halaman 59 (llm) ===

\section*{Library Crypto.Cipher}

\begin{itemize}
    \item Pustaka \texttt{Crypto.Cipher} adalah \textbf{package cipher} di dalam Python.
    \url{https://www.pycrypto.org/api/2.6/Crypto.Cipher-module.html}
    \item Di dalamnya terdapat beberapa modul beberapa cipher, salah satunya RC4 dengan nama modul \texttt{Crypto.Cipher.ARC4}
    \item Cara memanggil Pustaka \texttt{Crypto.Cipher}
    \begin{center}
        \texttt{>> from Crypto.Cipher import ARC4}
    \end{center}
    dengan asumsi sudah meng-install \textit{pycryptodome} dengan menggunakan perintah \texttt{pip}:
    \begin{center}
        \texttt{>> pip install pycryptodome}
    \end{center}
    \item Contoh program Python untuk enkripsi pesan teks dengan RC4 pada slide berikut
\end{itemize}

\end{document}
